% 图契约：不同形状的退磁场；图中的端口、几何、箭头和方程对应本题。
% 版面：沿用前三模块的几何图+推导分区；目标印宽17cm，简图至中等密度。
% 布局：左侧真实结构，右侧关键关系；真实电路边不作为架构流程箭头处理。
\documentclass[tikz,border=7pt]{standalone}
\usepackage{fontspec,xeCJK,amsmath,amssymb}
\IfFontExistsTF{Times New Roman}{\setmainfont{Times New Roman}}{\setmainfont{TeX Gyre Termes}}
\IfFontExistsTF{Heiti SC}{\setCJKmainfont[AutoFakeBold=1.5]{Heiti SC}}{\setCJKmainfont[AutoFakeBold=1.5]{FandolHei-Regular}}
\usetikzlibrary{arrows.meta,bending,calc,patterns,positioning,decorations.pathmorphing}
\definecolor{ink}{HTML}{243745}
\definecolor{blue}{HTML}{275DAB}
\definecolor{green}{HTML}{177C63}
\definecolor{red}{HTML}{BD4935}
\definecolor{orange}{HTML}{AA6C16}
\definecolor{purple}{HTML}{715393}
\tikzset{every node/.style={font=\fontsize{9.3}{12.5}\selectfont,text=ink},
arrow/.style={-{Stealth[length=5pt 1.5,width'=0pt 3]},shorten >=2pt,shorten <=1pt,line width=1pt},
arrow short/.style={-{Stealth[length=3pt 1.0,width'=0pt 2.5]},shorten >=1pt,shorten <=0pt,line width=.8pt},
axis/.style={arrow,draw=ink!60,line width=.6pt},
vec/.style={arrow,draw=blue},
wire/.style={draw=ink!75,line width=.8pt},
guide/.style={dashed,draw=ink!45,line width=.5pt},
card/.style={rounded corners=2pt,draw=blue!65,fill=blue!4,align=center,text width=6.5cm,inner xsep=9pt,inner ysep=7pt},
title/.style={font=\fontsize{10.5}{14}\selectfont\bfseries,text width=15cm,align=center},
label/.style={font=\fontsize{8.4}{11}\selectfont,align=center},
note/.style={label,text width=6.1cm}}
\begin{document}
\begin{tikzpicture}[x=1cm,y=1cm]
\path[use as bounding box] (-8,-3.5) rectangle(8,3.65);
\node[title] at (0,3.25) {理论图\quad 不同形状的退磁场};
\begin{scope}[shift={(-5,0)}]\draw[ink,fill=green!4](0,.6)circle(.9);\draw[vec,green](0,-.2)--(0,1.3)node[right]{$\mathbf M$};\node at(0,2.35){球体};\node at(0,-1.4){$N=1/3$};\end{scope}\begin{scope}[shift={(0,0)}]\draw[ink,fill=green!4](-.4,-.6)rectangle(.4,1.8);\draw[vec,green](0,-.2)--(0,1.3)node[right]{$\mathbf M$};\node at(0,2.35){细长圆柱};\node at(0,-1.4){$N\simeq0$};\end{scope}\begin{scope}[shift={(5,0)}]\draw[ink,fill=green!4](-1.5,.4)rectangle(1.5,.8);\draw[vec,green](0,-.2)--(0,1.3)node[right]{$\mathbf M$};\node at(0,2.35){法向磁化薄片};\node at(0,-1.4){$N\simeq1$};\end{scope}\node[card,text width=13cm]at(0,-2.6){$\mathbf H_d=-N\mathbf M$；退磁因子由形状和磁化方向共同决定};
\end{tikzpicture}\end{document}